\section{Introduction}

The purpose of vision is not exhausted by recognition. For humans, visual perception supports localization, action, memory, navigation, and imagination because it helps construct a structured model of the external world from partial retinal evidence. In psychophysics and systems neuroscience, this is precisely why researchers do not stop at isolated accuracy scores: they rely on controlled stimuli, repeated probing, perturbations, and cross-condition comparisons to infer what kind of representation an observer actually maintains \cite{kriegeskorte2008rsa,yamins2016dicarlo}. If we want to understand what a VLM really sees, we should adopt the same posture. A model should not be treated merely as a question-answering device; it should be treated as an observer whose behavior may or may not imply a coherent scene representation.

Current VLM evaluation falls short of this standard. Many benchmarks reduce spatial competence to isolated judgments such as ``left of,'' ``closer than,'' or ``in front of.'' A correct answer to one question, however, does not imply that the model has formed a global scene. Even many individually correct answers may be mutually incompatible once interpreted jointly. A model may appear competent at local distance comparison while failing to assemble those judgments into a single consistent layout; it may produce plausible directional answers by leaning on language priors or surface heuristics, yet those answers may immediately collapse when placed into a unified coordinate system. What most existing spatial benchmarks measure, therefore, is often local relational answering rather than reconstructable spatial belief.

Our motivation comes from a basic fact in ordinal embedding and scene reconstruction: when ordinal relations are sufficiently dense and sufficiently accurate, scene geometry can be approximately recovered up to similarity transforms \cite{agarwal2007gnmds,ma2018robustordinal,dokmanic2015edm}. This suggests a new way to evaluate VLMs. If we ask a model enough relative-distance and relative-direction questions, the resulting answer set becomes a form of behavioral tomography. QRR and TRR are not just tasks; they are probes of the model's externalized scene belief. QRR tests whether the model preserves the relative near-far structure among objects. TRR tests whether it maintains a coherent reference frame for direction. If the induced constraint set supports a reconstruction close to the ground-truth scene, then the model likely maintains a coherent scene-level representation. If not, then local success is better explained as a patchwork of priors, templates, or weak independent heuristics.

This is the motivation behind \textbf{Ordinary-Bench}. Rather than adding yet another spatial QA dataset, we treat a VLM as an observer to be behaviorally reverse-engineered. We generate controlled 3D-rendered scenes, probe them densely with QRR (quaternary relative relations) and TRR (ternary clock-direction relations), and then interpret the model's answers as ordinal and directional constraints whose consistency and reconstructability can be tested. The use of controlled synthetic scenes is deliberate. For the question ``what scene does the model actually externalize?'', exact geometric ground truth, controlled complexity, and repeatable perturbations are more valuable than raw ecological realism \cite{johnson2017clevr,yi2020clevrer}.

The empirical picture that motivates this paper is already clear. On the 540 scenes currently completed, single-view QRR reaches only about 70\% accuracy. More strikingly, adding multiple views does not improve this capability; instead, multi-view QRR degrades toward chance. TRR remains near chance throughout. This means that current VLMs may retain a weak and fragile notion of relative distance, but they still fail to form a reliable scene-level directional representation. Even more importantly, supervised fine-tuning or reinforcement learning on more than 200K relational QA pairs does not yield a meaningful improvement. This negative result matters: simply injecting more relational supervision into a model does not automatically produce a geometrically coherent scene belief. If evaluation stops at question-level accuracy, this structural failure is easy to miss.

From an application perspective, the stakes go beyond spatial QA. If a VLM can support stable scene reconstruction from ordinal relations, then it effectively exposes a vectorized scene representation that is only weakly coupled to any particular metric scale. Once the ordinal structure is stable, absolute scales can be aligned afterward using task-specific priors, empirical calibration, or ground-truth anchors. This creates a path toward scene reconstruction, localization, retrieval, and structured scene memory from relational queries alone. Ordinary-Bench is therefore not only about whether a model can answer spatial questions, but whether it has reached the more fundamental ability to recover a manipulable spatial structure from vision.

Our contributions are three-fold. First, we introduce a psychophysics-inspired evaluation framework for VLM spatial understanding, using controlled 3D scenes and dense QRR/TRR probing to study the model as a behaviorally observable observer. Second, we bring the idea of reconstructability from ordinal embedding into multimodal evaluation: enough relational answers should, in principle, reveal whether the model holds a coherent scene belief. Third, we present a set of already informative negative findings: limited single-view QRR success does not translate into stable spatial belief, multi-view input can collapse rather than improve performance, TRR remains near chance, and large-scale relational fine-tuning or RL does not substantially change this picture. Taken together, these results push the evaluation question from ``can a VLM answer spatial questions?'' to the stricter question: \emph{does the model actually construct a coherent, reconstructable world from visual input?}
